% =====================================================
% Eigenschaften des Elektrons: Masse, Ladung und Spin
% =====================================================
\subsection{Eigenschaften des Elektrons: Masse, Ladung und Spin}
\label{sec:eigenschaften-elektron}

Das Elektron\index{Elektron} ist eines der einfachsten bekannten Teilchen und zugleich
eines der folgenreichsten. Nach heutigem Wissen besitzt es keine erkennbare innere
Struktur. Es erscheint in der modernen Teilchenphysik\index{Teilchenphysik} als
elementares Teilchen\index{Elementarteilchen}. Dennoch besitzt es feste physikalische
Eigenschaften, durch die es sich eindeutig charakterisieren lässt: Masse\index{Masse},
elektrische Ladung\index{Ladung} und Spin\index{Spin}.

Die Masse des Elektrons ist außerordentlich klein. Sie beträgt ungefähr
\[
m_\mathrm{e} \approx 9{,}109 \cdot 10^{-31}\,\mathrm{kg}.
\]
Damit ist das Elektron viel leichter als ein Proton\index{Proton}. Gerade diese geringe
Masse ist entscheidend für viele Eigenschaften der Materie. Elektronen können sich in
Atomen\index{Atom}, Molekülen\index{Molekül} und Festkörpern\index{Festkörper}
vergleichsweise leicht umordnen. Deshalb bestimmen sie chemische Bindungen,
elektrische Leitfähigkeit und viele optische Eigenschaften von Stoffen.

Die zweite zentrale Eigenschaft ist die elektrische Ladung. Das Elektron trägt die
negative Elementarladung\index{Elementarladung}
\[
q_\mathrm{e} = -e
= -1{,}602\,176\,634 \cdot 10^{-19}\,\mathrm{C}.
\]
Das Minuszeichen ist dabei keine Nebensache. Es legt fest, wie das Elektron auf
elektrische Felder\index{elektrisches Feld} reagiert und wie es mit positiv geladenen
Teilchen wechselwirkt. Die elektrische Anziehung zwischen negativ geladenen Elektronen
und positiv geladenen Atomkernen\index{Atomkern} ist eine der Grundlagen für die
Stabilität der Atome.

Die dritte Eigenschaft ist der Spin. Das Elektron besitzt den Spin
\[
s = \frac{1}{2}.
\]
Der Spin ist eine rein quantenmechanische Eigenschaft. Man darf ihn nicht so verstehen,
als würde das Elektron wie eine kleine Kugel um die eigene Achse rotieren. Diese
Vorstellung ist anschaulich, aber irreführend. Der Spin verhält sich zwar in vielen
Rechnungen wie ein innerer Drehimpuls\index{Drehimpuls}, er ist jedoch kein klassischer
Rotationszustand eines ausgedehnten Körpers.

Mit dem Spin ist auch ein magnetisches Moment\index{magnetisches Moment} verbunden.
Deshalb reagiert das Elektron auf Magnetfelder\index{Magnetfeld}. Diese Eigenschaft
spielt eine wichtige Rolle in der Atomphysik\index{Atomphysik}, in der Chemie, in der
Festkörperphysik\index{Festkörperphysik} und in modernen Messverfahren wie der
Magnetresonanztomographie\index{Magnetresonanztomographie}.

Masse, Ladung und Spin reichen aus, um das Elektron als Teilchen eindeutig zu
charakterisieren. Zugleich zeigen sie bereits, warum das Elektron so wichtig ist:
Seine Masse macht es beweglich, seine Ladung macht es elektromagnetisch wirksam,
und sein Spin führt zu quantenmechanischen und magnetischen Effekten. Aus diesen
wenigen Eigenschaften entsteht ein großer Teil der physikalischen Vielfalt der
sichtbaren Welt.

\begin{PhysicsBox}[Die drei Grundeigenschaften des Elektrons]
	\label{box:grundeigenschaften}
	Das Elektron wird durch drei grundlegende Eigenschaften charakterisiert:
	\[
	m_\mathrm{e} \approx 9{,}109 \cdot 10^{-31}\,\mathrm{kg},
	\quad
	q_\mathrm{e} = -1{,}602\,176\,634 \cdot 10^{-19}\,\mathrm{C},
	\quad
	s = \frac{1}{2}.
	\]
	
	Die Masse beschreibt seine Trägheit, die Ladung seine elektromagnetische
	Wechselwirkung und der Spin seine quantenmechanische Drehimpulseigenschaft.
	Gerade die Kombination dieser drei Eigenschaften macht das Elektron zu einem
	zentralen Baustein der Materie.
\end{PhysicsBox}